\chapter{Revisão Sistemática da Literatura}
\label{cap:revisao_sistematica}

% -----------------------------------------------------------------
% Função do capítulo: revisar de forma sistemática, dentro da
% literatura de Computação, os métodos orientados a dados que
% endereçam o problema descrito no Capítulo 2. A revisão segue o
% protocolo de Kitchenham e cobre as duas famílias de técnicas que
% sustentam a proposta (detecção de anomalias em séries operacionais
% e estratégias de integração com explicabilidade), além de fechar
% com a lacuna que o trabalho preenche.
%
%   3.1  Protocolo da revisão (Kitchenham): pergunta de pesquisa
%        orientadora, PICOC, string de busca, bases, critérios de
%        inclusão e exclusão, fluxograma PRISMA.
%   3.2  Detecção de anomalias em dados industriais. RQ1.
%   3.3  Sistemas híbridos e explicabilidade. RQ2 e RQ4.
%   3.4  Síntese da revisão e identificação da lacuna que motiva
%        a proposta do Capítulo 4.
%
% Sementes para bootstrap:
%   Chandola et al. (2009), Pang et al. (2021), Liu et al. (2008).
% Protocolo detalhado em:
%   docs/dissertacao/protocolo_revisao.md
% -----------------------------------------------------------------

Este capítulo conduz a revisão sistemática da literatura de Computação aplicada a sistemas industriais, com foco nos métodos orientados a dados que se propõem a detectar anomalias em séries operacionais e a integrar evidências de múltiplas fontes em diagnósticos interpretáveis. A revisão segue o protocolo de Kitchenham e Charters, com pergunta de pesquisa orientadora, estratégia de busca em bases científicas reconhecidas, critérios formais de inclusão e exclusão e síntese dos trabalhos selecionados. As tradições de diagnóstico oriundas do próprio domínio industrial, baseadas em ferramentas de qualidade e em análise estatística clássica de confiabilidade, foram apresentadas no Capítulo~\ref{cap:fundamentacao} e não constituem objeto desta revisão sistemática.


\section{Protocolo da revisão}
\label{sec:protocolo}

% Função: detalhar o protocolo Kitchenham aplicado neste trabalho.
% Conteúdo a desenvolver:
%   - Pergunta de pesquisa orientadora da revisão.
%   - Estratégia PICOC (População, Intervenção, Comparação, Outcome, Contexto).
%   - String de busca aplicada nas bases.
%   - Bases consultadas (Scopus, IEEE Xplore, ACM Digital Library, Web of Science).
%   - Critérios de inclusão e exclusão.
%   - Fluxograma PRISMA com contagem de artigos por etapa.
%   - Avaliação de qualidade dos artigos selecionados.
%   - Síntese e extração de dados dos artigos.
% Protocolo detalhado em docs/dissertacao/protocolo_revisao.md


\section{Detecção de anomalias em dados industriais}
\label{sec:anomalias}

A detecção de anomalias compreende o conjunto de técnicas dedicadas a identificar observações que se afastam do padrão esperado em um conjunto de dados. A literatura distingue três tipos canônicos. Anomalias pontuais correspondem a observações isoladas que diferem do restante do conjunto. Anomalias contextuais são observações que só se mostram anômalas quando avaliadas em relação ao seu contexto, como uma leitura de temperatura normal no inverno mas atípica no verão. Anomalias coletivas, por sua vez, são grupos de observações que formam um padrão atípico mesmo que cada elemento isolado pareça normal \cite{chandola2009anomaly}. A escolha do método depende também do regime de aprendizado disponível: técnicas supervisionadas exigem rótulos de normalidade e anomalia para cada amostra de treino, técnicas semi-supervisionadas exigem apenas exemplos da classe normal, e técnicas não supervisionadas operam sem rótulo algum, identificando padrões raros ou pouco frequentes a partir da estrutura intrínseca dos dados. No cenário desta dissertação, em que rompimentos são registrados sem atribuição de causa, o regime não supervisionado é o único viável e orienta toda a discussão da seção.

Quando o objeto de detecção é uma linha de produção, a tarefa adquire complicadores específicos. Os sensores registram dezenas de variáveis simultaneamente, com correlações estruturais entre si que refletem o acoplamento físico do processo, como a relação direta entre a temperatura do forno e a velocidade da linha. O regime de operação não é estacionário: os valores de referência operacional mudam ao longo do dia, materiais distintos são processados em sequência e a linha passa por transientes legítimos como rampas de partida e trocas de produto, que precisam ser distinguidos de anomalias reais. Além disso, o conceito de comportamento normal evolui ao longo do tempo, fenômeno conhecido como deriva conceitual (\emph{concept drift}), o que exige que o modelo seja capaz de se adaptar sem comprometer a sensibilidade \textcolor{red}{[ref: revisão sobre concept drift, p. ex. Gama et al. 2014 ou Lu et al. 2018]}. Esses fatores tornam abordagens univariadas e estacionárias, como as cartas de controle clássicas discutidas na Seção~\ref{sec:diagnostico}, insuficientes para o tipo de dado disponível.

As técnicas de detecção de anomalias se organizam em famílias distintas conforme o princípio de discriminação que adotam. Métodos estatísticos clássicos, como a estatística $T^2$ de Hotelling apresentada na Seção~\ref{sec:diagnostico}, modelam a distribuição esperada dos dados e sinalizam como anômalas as observações cuja verossimilhança fica abaixo de um limiar pré-estabelecido. Métodos baseados em densidade, como o \emph{Local Outlier Factor} (LOF), comparam a densidade local de cada ponto com a densidade de seus vizinhos no espaço de atributos, sinalizando como anômalas as observações que residem em regiões esparsas \textcolor{red}{[ref: Breunig et al. 2000 (\emph{LOF: Identifying Density-Based Local Outliers}) e Aggarwal 2017 (\emph{Outlier Analysis})]}. Métodos baseados em isolamento exploram a observação de que pontos anômalos tendem a ser isolados em poucos cortes aleatórios do espaço, dispensando estimativas explícitas de densidade ou distribuição. Modelos profundos, como autoencoders e redes recorrentes do tipo \emph{Long Short-Term Memory} (LSTM), oferecem alternativas para situações em que padrões complexos e relações temporais de longo prazo precisam ser capturados, ao custo de maior demanda computacional e de dados de treino \textcolor{red}{[ref: Pang et al. 2021 (\emph{Deep Learning for Anomaly Detection: A Review}), ACM Computing Surveys]}. A escolha entre famílias depende de fatores como dimensionalidade, volume de dados, requisitos de interpretabilidade e custo computacional.

Dentre as famílias apresentadas, o método de isolamento conhecido como \emph{Isolation Forest} é o adotado nesta dissertação \textcolor{red}{[ref: Liu, Ting e Zhou 2008 (\emph{Isolation Forest}), ICDM, ou Liu et al. 2012 (\emph{Isolation-Based Anomaly Detection}), TKDD]}. O algoritmo opera construindo um \emph{ensemble} de árvores binárias em que cada divisão é feita selecionando aleatoriamente um atributo e um valor de corte dentro do intervalo observado desse atributo. Sob essa partição aleatória, pontos anômalos tendem a ser separados do conjunto principal em poucos níveis da árvore, enquanto pontos normais exigem partições sucessivas até serem isolados. O escore de anomalia atribuído a cada observação deriva da profundidade média de isolamento ao longo das árvores do ensemble, normalizado para o intervalo entre zero e um, com valores próximos de um indicando alta probabilidade de anomalia. O método se caracteriza pela escalabilidade em alta dimensionalidade, pela ausência de pressupostos paramétricos sobre a distribuição dos dados e pelo custo computacional linear no tamanho do conjunto. Os principais hiperparâmetros são o número de árvores do ensemble, o tamanho da sub-amostra usada para construir cada árvore e o parâmetro de contaminação esperada, que define o limiar de classificação entre escores normais e anômalos.

A avaliação de detectores de anomalia treinados sem rótulo apresenta dificuldade própria. Métricas clássicas como precisão, \emph{recall} e área sob a curva ROC (\emph{Receiver Operating Characteristic}) dependem da existência de uma verdade de referência (\emph{ground truth}) que indique quais observações são efetivamente anômalas, condição não satisfeita no cenário aqui considerado. A literatura propõe estratégias alternativas para lidar com essa limitação. Uma é a avaliação por especialistas, que examinam manualmente um subconjunto das observações sinalizadas pelo modelo e julgam se o sinal é coerente com o comportamento esperado da operação. Outra é a comparação com eventos conhecidos, na qual os escores produzidos pelo modelo são confrontados com registros de incidentes documentados a posteriori. A análise de consistência interna verifica se observações próximas no tempo ou no espaço de atributos recebem escores compatíveis, sob a hipótese de que a transição entre estados operacionais é gradual. Estudos de robustez, por fim, treinam e testam o modelo com diferentes recortes e parametrizações para avaliar a estabilidade dos resultados frente a variações de calibragem \textcolor{red}{[ref: Goldstein e Uchida 2016 (\emph{A Comparative Evaluation of Unsupervised Anomaly Detection Algorithms}) e Pang et al. 2021]}. Essas estratégias não substituem rótulos, mas oferecem evidências indiretas sobre o comportamento do modelo.

A literatura registra aplicações de detecção de anomalia em diversos domínios de manufatura, incluindo monitoramento preditivo de máquinas, controle de qualidade de produto e identificação de desvios de processo \textcolor{red}{[ref: surveys de manutenção preditiva e Indústria 4.0, p. ex. Carvalho et al. 2019 ou Zhao et al. 2019]}. No contexto da fabricação de fios esmaltados, contudo, técnicas isoladas de anomalia apresentam limite estrutural: o modelo é capaz de sinalizar que o comportamento operacional da máquina se afastou do esperado em determinado intervalo, mas é incapaz, por construção, de atribuir esse desvio a uma causa específica. Quando o rompimento decorre de defeito de matéria-prima formado antes do início da linha analisada, conforme discutido na Seção~\ref{sec:problema_rompimentos}, os sensores podem nem registrar comportamento atípico, pois a falha viaja embutida no fio e se manifesta apenas em um ponto físico distante de sua origem. Além disso, a detecção de anomalias por séries operacionais ignora a evidência agregada presente no histórico estatístico de rompimentos por lote, material, fornecedor e máquina, apresentado na Seção~\ref{sec:estatistica_historica} do Capítulo~\ref{cap:fundamentacao}. A integração entre essas duas fontes de evidência, por meio de arquiteturas híbridas e interpretáveis, é objeto da Seção~\ref{sec:hibridos_xai}.


\section{Sistemas híbridos e explicabilidade}
\label{sec:hibridos_xai}

% Função retórica: mostrar que combinar a detecção de anomalias em
% séries operacionais com a análise estatística histórica de falhas
% é a direção promissora, mas pouco consolidada no domínio de fios
% esmaltados, e apontar a lacuna.
% RQ2 e RQ4.
% Conteúdo a desenvolver:
%   - Estratégias de combinação: regras, ensemble, fusão de
%     evidências, sistemas em camadas.
%   - Importância de resultados interpretáveis em chão de fábrica.
%   - Interpretabilidade aplicada a manutenção, qualidade e
%     engenharia.
% Conclui com o parágrafo da lacuna:
%   - Existe maturidade nas duas famílias isoladas.
%   - Há trabalhos pontuais de integração em outros domínios.
%   - Falta uma arquitetura híbrida e explicável aplicada ao
%     domínio específico de fios esmaltados, com regra de
%     cruzamento entre máquinas para diferenciar falha de
%     equipamento de suspeita de matéria-prima.
%   - Essa lacuna motiva a proposta do Capítulo~\ref{cap:proposta}.


\section{Síntese da revisão e identificação da lacuna}
\label{sec:sintese_revisao}

% Função: consolidar achados da revisão, organizar os trabalhos
% selecionados em tabela de síntese e amarrar a lacuna que a proposta
% do Capítulo 4 endereça.
% Conteúdo a desenvolver:
%   - Tabela com os trabalhos mais relevantes, classificados por
%     família (anomalia, estatística histórica, híbrido) e por
%     domínio de aplicação.
%   - Pontos fortes e limites de cada família.
%   - Identificação explícita da lacuna que o trabalho preenche.
%   - Como a Proposta do Capítulo~\ref{cap:proposta} se posiciona em
%     relação aos trabalhos selecionados.
